% Part: sets-functions-relations
% Chapter: sets
% Section: pairs-and-products

\documentclass[../../../include/open-logic-section]{subfiles}

\begin{document}

\olfileid{sfr}{set}{pai}
\olsection{Pasangan, Tupel, Produk Kartesius}

\begin{explain}
Saka ekstensionalitas bisa didudut yen anggota himpunan ora
duwe urutan. Mula, yen arep nggambarake urutan, kita nganggo
\emph{pasangan urut} $\tuple{x, y}$. Ing pasangan tanpa urutan
$\{x, y\}$, urutan ora wigati: $\{x, y\} = \{y, x\}$.
Ing pasangan urut, urutan iku wigati: yen $x
\neq y$, mula $\tuple{x, y} \neq \tuple{y, x}$.

Kepriye kudune kita ngerteni pasangan urut ing teori himpunan?
Kang wigati, kita arep njaga gagasan yen pasangan urut iku padha
yen lan mung yen anggota kapisane padha lan anggota kapindhone
uga padha, yaiku:
\[
  \tuple{a, b}= \tuple{c, d}\text{ yen lan mung yen }a = c \text{ lan }b=d.
\]
Pasangan urut bisa ditemtokake ing teori himpunan nganggo
definisi Wiener--Kuratowski.
\end{explain}

\begin{defn}[Pasangan urut]\ollabel{wienerkuratowski}
	$\tuple{a, b} = \{\{a\}, \{a, b\}\}$.
\end{defn}

\begin{prob}
	Kanthi nggunakake \olref[sfr][set][pai]{wienerkuratowski},
	buktekna yen $\tuple{a,b}= \tuple{c, d}$ yen lan mung yen
	$a = c$ lan $b=d$.
\end{prob}

\begin{explain}
Sawise netepake definisi pasangan urut, kita bisa nggunakake
kanggo nemtokake himpunan liyane. Tuladhane, kadhangkala kita
uga mbutuhake urutan luwih saka rong objek, kayata
\emph{tupel telu} $\tuple{x, y, z}$,
\emph{tupel papat} $\tuple{x, y, z, u}$, lan sateruse.
Tupel telu bisa dianggep pasangan urut khusus kang anggota
kapisane dhewe arupa pasangan urut:
$\tuple{x, y, z}$ yaiku $\tuple{\tuple{x, y},z}$.
Semono uga kanggo tupel papat: $\tuple{x,y,z,u}$ yaiku
$\tuple{\tuple{\tuple{x,y},z},u}$, lan sateruse.
Kanthi umum, kita nyebut \emph{tupel-$n$ urut}
$\tuple{x_1, \dots, x_n}$.

Sawetara himpunan pasangan urut, utawa tupel-$n$ urut liyane,
bakal migunani.
\end{explain}

\begin{defn}[Produk Kartesius]
Kanggo himpunan $A$ lan $B$, \emph{produk Kartesius}
$A \times B$ ditemtokake kanthi
\[
  A \times B = \Setabs{\tuple{x, y}}{x \in A \text{ lan } y \in B}.
\]
\end{defn}

\begin{ex}
Yen $A = \{0, 1\}$ lan $B = \{1, a, b\}$, mula produke yaiku
\[
A \times B = \{ \tuple{0, 1}, \tuple{0, a}, \tuple{0, b},
    \tuple{1, 1}, \tuple{1, a}, \tuple{1, b} \}.
\]
\end{ex}

\begin{ex}
Yen $A$ iku himpunan, produk $A$ karo awake dhewe,
$A \times A$, uga ditulis~$A^2$. Iki himpunan
\emph{kabeh} pasangan $\tuple{x, y}$ kanthi $x, y \in A$.
Himpunan kabeh tupel telu $\tuple{x, y, z}$ yaiku $A^3$,
lan sateruse. Kita bisa menehi definisi rekursif:
\begin{align*}
  A^1 & = A\\
  A^{k+1} & = A^k \times A
\end{align*}
\end{ex}

\begin{prob}
Dhaptarna kabeh !!{element}s saka $\{1, 2, 3\}^3$.
\end{prob}

\begin{prop}\ollabel{cardnmprod}
Yen $A$ duwe $n$ !!{element}s lan $B$ duwe $m$ !!{element}s,
mula $A\times B$ duwe $n\cdot m$ anggota.
\end{prop}

\begin{proof}
Kanggo saben !!{element}~$x$ ing~$A$, ana $m$ !!{element}s
awujud $\tuple{x, y} \in A \times B$.
Tetepna $B_x = \Setabs{\tuple{x, y}}{y\in B}$.
Amarga saben $x_1 \neq x_2$ ndadekake
$\tuple{x_1, y} \neq\tuple{x_2, y}$, mula
$B_{x_1} \cap B_{x_2} = \emptyset$.
Nanging yen $A = \{x_1,\dots, x_n\}$, mula
$A \times B = B_{x_1} \cup \dots \cup B_{x_n}$,
saengga duwe $n\cdot m$ !!{element}s.

Kanggo nggambarake iki, tatanen !!{element}s saka~$A \times B$
ing larik lan kolom:
\[
\begin{array}{rcccc}
  B_{x_1} = & \{\tuple{x_1, y_1} & \tuple{x_1, y_2} & \dots & \tuple{x_1, y_m}\}\\
  B_{x_2} = & \{\tuple{x_2, y_1} & \tuple{x_2, y_2} & \dots & \tuple{x_2, y_m}\}\\
  \vdots & & \vdots\\
  B_{x_n} = & \{\tuple{x_n, y_1} & \tuple{x_n, y_2} & \dots & \tuple{x_n, y_m}\}
\end{array}
\]
Amarga kabeh $x_i$ beda-beda lan kabeh $y_j$ beda-beda,
ora ana rong pasangan ing tatanan iki kang padha,
lan cacahe kabeh pasangan yaiku $n\cdot m$.
\end{proof}

\begin{prob}
Tuduhna, kanthi induksi tumrap~$k$, yen kanggo kabeh $k \ge 1$,
yen $A$ duwe $n$ !!{element}s, mula $A^k$ duwe
$n^k$ !!{element}s.
\end{prob}

\begin{ex}
Yen $A$ iku himpunan, \emph{tembung} ing~$A$ yaiku saben
urutan !!{element}s saka~$A$. Urutan bisa dianggep
tupel-$n$ kang ngemot !!{element}s saka~$A$.
Tuladhane, yen $A = \{a, b, c\}$, urutan "$bac$" bisa
dianggep minangka tupel telu~$\tuple{b, a, c}$.
Tembung, yaiku urutan lambang, wigati banget ing ilmu komputer.
Miturut konvensi, !!{element}s saka~$A$ dianggep urutan
dawane~$1$, lan $\emptyset$ minangka urutan dawane~$0$.
Mula himpunan \emph{kabeh} tembung ing~$A$ yaiku
\[
A^* = \{\emptyset\} \cup A \cup A^2 \cup A^3 \cup \dots
\]
\end{ex}

\end{document}
